\section{Scope and limitations}\label{sec:scope}

\subsection{Summary}\label{sec:scope:nonclaims}

This section collects what the paper does not claim.
\Cref{tab:nonclaims} summarizes it, and the subsections below give the
details; they govern how the earlier sections should be read.

\begin{table}[H]
  \centering
  \caption{What the paper does not claim.}
  \label{tab:nonclaims}
  \small
  \renewcommand{\arraystretch}{1.1}
  \begin{tabularx}{\textwidth}{@{}>{\raggedright\arraybackslash}p{0.30\textwidth}>{\raggedright\arraybackslash}X@{}}
    \toprule
    Topic & Limitation \\
    \midrule
    Birch--Swinnerton-Dyer & Outside the scope of this paper; not used as evidence. \\
    Hodge & A prediction only (\Cref{sec:predictions:hodge}). \\
    Unconditional results & None.  Each application keeps its own premises; the Gaussian constructions supply none of them. \\
    Recognition sources & Named premises, not derived; their acceptance and any admission laws are hypotheses. \\
    Recognition-mode schemas & Conditional schemas illustrated by two applications; not proved. \\
    Non-descending translation & Not stated or used as a theorem; represented only by the calibration family and the open gates. \\
    \bottomrule
  \end{tabularx}
\end{table}


\subsection{The Birch--Swinnerton-Dyer conjecture}\label{sec:scope:bsd-out}

The paper treats three applications: Navier--Stokes regularity, the
Riemann hypothesis, and \(P\) versus \(NP\).  A companion paper treats
the Birch--Swinnerton-Dyer conjecture \cite{TsiokosBSD2026,Wiles2000BSD};
it is not analysed here and is not used as evidence.  This is a choice
of scope, not a statement about that paper.

\subsection{The Hodge conjecture}\label{sec:scope:hodge-prediction-only}

No formed closure for the Hodge conjecture exists in this series.  The
Hodge conjecture appears only in the prediction of
\Cref{sec:predictions:hodge}, which proposes it as a third test case.
The paper does not claim that the conjecture holds, that it satisfies
the recognition-mode template, or that a Hodge recognition source has
been found.

\subsection{No unconditional claim}\label{sec:scope:standard-zfc}

Every application cited here is conditional on its own premises, and
the framework account also uses the SB closure assumption.  This paper
strengthens none of the applications to an unconditional theorem and
does not describe any of them as unconditional.  The premises are, for
Navier--Stokes, a spectral closure, or a physical-window closure
together with an assumed terminal-divergence input, with initial-data
hypotheses; for the Riemann hypothesis, a supplied
Selberg recognition source and an identification with the classical
zeros; and for \(P\) versus \(NP\), that a fixed tagged witness
selector is not computable in polynomial time, or, in the family form,
a named negative source together with the admission of search events.

\Cref{thm:joint-gaussian-return} is unconditional as an identity on the
objects it names.  Its zero-side premise \(\mathcal A_{a,\nu,t}\) is
equivalent to the Riemann hypothesis, and its Navier--Stokes conclusions
assume a separate closure hypothesis; the identity establishes neither
and does not combine them into one closure law.
\Cref{thm:concrete-triad-gaussian-xi} likewise proves exact identities.
Its SAT sum has \(2^{n}\) terms, and its positivity gives no
polynomial-time procedure.  It transfers no premise from one
application to another, and the impossibility statement of
\Cref{sec:triad-bridge} excludes only maps that preserve the probe and
send the reflected displacement to a satisfiable SAT state.

\subsection{Recognition sources are premises}\label{sec:scope:recognition-source-imports}

The recognition sources are named premises, not derived results.  The
P-versus-NP source is recorded in that paper; its acceptance and the
admission of search events are separate premises.  The Selberg source
is recorded in the Riemann-hypothesis paper.
The paper never treats these sources as derived from the framework or
produced by the recognition-mode template.

\subsection{The recognition-mode schemas are not proved}\label{sec:scope:universal-lifting-deferred}

The four statements of \Cref{sec:recognition} are conditional schemas,
illustrated by the two companion applications and not proved.
Turning them into theorems needs definitions of their predicates and
the arguments listed in \Cref{sec:predictions:theorem-grade-lifting}.  The relation between
\Cref{thm:attack-foreclosure-v5} and \Cref{thm:attack-foreclosure-v4}
is unchanged by this: the former restricts the latter to derivation
routes and adds the recognition-mode clauses.

\subsection{Non-descending translation}
\label{sec:scope:m6-theorem-7-folded}

No non-descending translation theorem for the P-versus-NP application
is stated or used here.  The transport content such a statement would
carry is represented in this paper only by the calibration family of
\Cref{sec:calibration} and the open gates of
\Cref{sec:predictions:functorial-gates}.
